%
% dial_settlement.tex
%
% A Z specification of how a dial is ENDED and SETTLED: the reader slot, the
% listener's start and stop, the inline reader's last check and release, the
% no-reader settle, and the Hub's own end.  It refines the two abstractions that
% docs/crash_attribution.tex makes (a "reader slot" that someone holds, a
% "settle" that someone performs) to the granularity of the code, because three
% rounds of the final review on the lux-ghhf branch found a path on which an
% ended dial is never settled or is settled undrained (R7-F1, R8 R1, and the
% three defects of round 9).  This is the recurrence rule: the same class of
% defect in two or more rounds is a state-space problem, and the state space is
% finite.
%
% THIS FILE IS A FIDELITY CONTROL: it is dial_settlement.tex with StRel0 changed,
% and reproduces the round-9 silent-failure trace (S1 violated: the wedge).  The
% three controls are:
%   docs/dial_settlement_start_fail_buggy.tex       a start that fails keeps the slot
%   docs/dial_settlement_stop_skip_drain_buggy.tex  the reader's exit reads ended twice
%   docs/dial_settlement_inline_strand_buggy.tex    a refused settle is dropped
% The code it governs: src/punt_lux/domain/hub/{display_connection,dial_reader,
% display_listener,listener_thread,display_replies,reply_sink,dial_generation,
% reader_slot}.py.  Companions: crash_attribution.tex (what a capture charges),
% listener_locks.tex (the lock order rq < sl < st < lk < qu).
%
% Type-check:  fuzz docs/dial_settlement.tex
% Model-check: probcli docs/dial_settlement.tex -model_check -p MAX_OPERATIONS 100000
%              probcli docs/dial_settlement.tex -model_check -nodead -bf \
%                -goal "<S1: wedge, or S2 to S5: a flag is set>"   (see Verification)
%
\documentclass[a4paper,10pt,fleqn]{article}
\usepackage[margin=1in]{geometry}
\usepackage{fuzz}

\begin{document}

\title{Dial Settlement: a Z Specification}
\author{Who ends a dial, who settles it, and who may let go}
\date{October 2026}
\maketitle

\section{What is modelled}

A \emph{dial} is one generation of the Hub's connection to the display.  It is
\emph{published} (current) from its birth until some one party \emph{captures}
it: takes its ledger, queues its death or carries its suspects, and makes it
not current.  Between the first end and the capture the dial is ended and
unsettled.  Four properties are claimed of that interval.

\begin{description}
  \item[S1] No dial stays ended (or end-requested), current and unsettled for
    ever: an end is always followed by a settle, by a party that was already
    running when the end was made.  A dial that is wedged is one that
    redial cannot replace, because it is still current.
  \item[S2] A display or unresponsive end is captured only after the socket has
    been drained (T14, O2 of the coverage audit).
  \item[S3] Exactly one capture per dial.
  \item[S4] The cause that is captured is the first one recorded.
  \item[S5] A display that has already died is never recorded as a Hub end: the
    cause that wins is the display's, whatever the Hub did next.
\end{description}

S1 is stated as a reachability check, not as a state invariant: ProB treats the
primed state schema as part of each operation's guard, so a violation written as
a conjunct disables the operation instead of reporting it.  S2 to S5 are flags
set by the operation that would violate them; the goal is a state with a flag
set.

\section{The design in one rule}

The code has three separate rules where one will do.  Settlement belongs to the
\emph{holder of the reader slot}; everyone else only \emph{asks}.

\begin{enumerate}
  \item \textbf{Capture authority is the slot.}  Only a holder of the slot
    (the listener thread, an inline reader, or a releaser taking over) captures a
    dial.  No other path calls the capture.  In code: \texttt{close\_if\_current}
    leaves the connection's interface; the slot hands its holder a
    \texttt{Settler}, the one object with a \texttt{settle()} method, and the
    listener, the inline reader and the no-reader settle are the only holders.
    A flag with a default (\texttt{by\_reader=True}) is gone, and with it the
    Hub's capture door.
  \item \textbf{Ends record or request; they never capture.}  A display or
    unresponsive end records its cause, first wins.  A Hub end sets a request
    flag (\texttt{hubreq}); it records no cause.  Both then \emph{ask for a
    settle}: take the slot if it is free and settle at once; if it is held, leave
    a debt on the slot.
  \item \textbf{Whoever lets go of the slot honours the debt.}  A release
    (listener start failure, inline exit) that finds a debt keeps the slot,
    settles, and looks again.  Every path that gives up the reader role releases
    the slot this way, including a start that raises because the dial was ended.
  \item \textbf{Settle is one idempotent step that reads the cause once.}  It
    drains the socket to the point where it would block, then looks at the
    outcome: end of file is the display's end, else a pending Hub request is the
    Hub's end, else (if an end was already recorded) that one stands.  The
    cause is read once, by this step, after the drain, so no end landing later
    can change what is captured or whether it was drained.
\end{enumerate}

Rules 2 and 4 together are what close S5: the Hub's end is decided at the
instant the socket is observed, and a display that died before that instant is
seen as dead.  The model assumes the observation is one linearization point
(the drain to \texttt{EAGAIN}, with the end-of-file test).  A display that dies
after it dies after the Hub's end; that is not a defect, and no rule can
distinguish it.

\section{Carriers and bounds}

One dial.  Every operation in the code is scoped to a generation: a reader, a
slot, a write failure, and a capture all name the dial they act on, and the Hub
redials only when the dial is no longer current.  A second dial therefore runs
in a disjoint product with the first, and cannot interact with it; the model
includes \texttt{Redial}, enabled only on an unpublished dial, to say so.  Five
actors, each of which acts at most once: the starter (\texttt{start\_on}), the
stopper (a Hub end: \texttt{stop} or \texttt{close}), a discoverer of a display
or unresponsive end (a write failure or the keepalive), one inline reader, and
the listener thread.  Socket content is zero or one unread acknowledgement,
which is enough to tell a drained capture from an undrained one.  The dial is
born holding the listener's slot (the registry's \texttt{open\_listening}) or
free (a library link).

\begin{zed}
  CAUSE ::= nocause | cdisp | cunresp | chub
\end{zed}

\begin{zed}
  SLOT ::= sfree | slistener | sinline
\end{zed}

\begin{zed}
  FLAG ::= fclear | fset
\end{zed}

\begin{zed}
  ACTOR ::= aSt | aSp | aDi | aIn | aRd
\end{zed}

The program counters.  \texttt{pst0}: the starter has not yet started the
listener.  \texttt{prel0}: it failed and owes a release.  \texttt{phub}: a Hub
end is to be requested.  \texttt{pend}: a display end is to be recorded.
\texttt{psu}: the no-reader settle is to be tried.  \texttt{pset}: this actor
holds the slot and owes a settle.  \texttt{prel}: it is about to release.
\texttt{iread}, \texttt{ilast}: the inline reader is reading, and has finished
reading and will make its last check.  \texttt{rtop}: the listener is at the top
of its loop.  \texttt{rnone}: no listener exists.

\begin{zed}
  PC ::= pidle | pst0 | prel0 | phub | pend | psu | pset | prel | pdone |
    iread | ilast | rtop | rnone | rdone
\end{zed}

\begin{schema}{Dial}
  ended : CAUSE \\
  slot : SLOT \\
  debt : FLAG \\
  current : FLAG \\
  sock : 0 \upto 1 \\
  capN : 0 \upto 2 \\
  hubreq : FLAG
\end{schema}

\begin{schema}{Env}
  launched : FLAG \\
  eof : FLAG \\
  redialed : FLAG
\end{schema}

\begin{schema}{Flags}
  s2bad : FLAG \\
  s3bad : FLAG \\
  s4bad : FLAG \\
  s5bad : FLAG
\end{schema}

\begin{schema}{Pc}
  pc : ACTOR \fun PC
\end{schema}

\begin{schema}{Dc}
  dcs : CAUSE
\end{schema}

\begin{schema}{RLoc}
  rcause : CAUSE \\
  rdec : FLAG
\end{schema}

\texttt{ended} is the first recorded cause.  \texttt{hubreq} is a request, not a
cause.  \texttt{eof} says that the display has died and its end of file is
queued behind whatever is unread; \texttt{sock} is that unread
acknowledgement.  \texttt{dcs} is the discoverer's cause, display or
unresponsive.  \texttt{rcause} and \texttt{rdec} are locals of the listener
that only the stop-skip-drain control uses.

\begin{schema}{Sys}
  Dial \\ Env \\ Flags \\ Pc \\ Dc \\ RLoc
\end{schema}

\begin{schema}{Init}
  Sys'
\where
  ended' = nocause \\
  debt' = fclear \\
  current' = fset \\
  sock' \in 0 \upto 1 \\
  capN' = 0 \\
  hubreq' = fclear \\
  launched' = fclear \\
  eof' = fclear \\
  redialed' = fclear \\
  s2bad' = fclear \\
  s3bad' = fclear \\
  s4bad' = fclear \\
  s5bad' = fclear \\
  dcs' = cdisp \\
  rcause' = nocause \\
  rdec' = fclear \\
  (slot' = slistener \land pc' = \{ aSt \mapsto pst0, aSp \mapsto pidle,
    aDi \mapsto pidle, aIn \mapsto pidle, aRd \mapsto rnone \}) \lor
  (slot' = sfree \land pc' = \{ aSt \mapsto pdone, aSp \mapsto pidle,
    aDi \mapsto pidle, aIn \mapsto pidle, aRd \mapsto rnone \})
\end{schema}

\section{The environment}

\begin{schema}{Die}
  \Xi Dial \\ \Delta Env \\ \Xi Flags \\ \Xi Pc \\ \Xi Dc \\ \Xi RLoc
\where
  eof = fclear \\
  eof' = fset \\
  launched' = launched \\
  redialed' = redialed
\end{schema}

The display dies.  Its end of file is queued behind any acknowledgement it has
already sent.

\begin{schema}{Redial}
  \Xi Dial \\ \Delta Env \\ \Xi Flags \\ \Xi Pc \\ \Xi Dc \\ \Xi RLoc
\where
  current = fclear \\
  redialed = fclear \\
  redialed' = fset \\
  launched' = launched \\
  eof' = eof
\end{schema}

The Hub dials again, and only on a dial that is no longer current.  A dial
that is ended but still current is never replaced: that is the wedge.

\begin{schema}{Halt}
  \Xi Sys
\where
  current = fclear \lor (ended = nocause \land hubreq = fclear)
\end{schema}

The model may stop only when the dial is settled or has not been ended.  A
state with the dial ended or end-requested, still current, and nothing enabled
is a deadlock, and probcli reports it: that is S1.

\section{The starter}

\begin{schema}{StOk}
  \Xi Dial \\ \Delta Env \\ \Xi Flags \\ \Delta Pc \\ \Xi Dc \\ \Xi RLoc
\where
  pc~aSt = pst0 \\
  current = fset \\
  ended = nocause \\
  hubreq = fclear \\
  launched' = fset \\
  eof' = eof \\
  redialed' = redialed \\
  pc' = pc \oplus \{ aSt \mapsto pdone, aRd \mapsto rtop \}
\end{schema}

\texttt{start\_on}: the dial is not closed, and the thread starts.

\begin{schema}{StFail}
  \Xi Dial \\ \Xi Env \\ \Xi Flags \\ \Delta Pc \\ \Xi Dc \\ \Xi RLoc
\where
  pc~aSt = pst0 \\
  pc' = pc \oplus \{ aSt \mapsto prel0 \}
\end{schema}

\texttt{start\_on} fails: the dial was found ended or requested, or the thread
could not be started.  Nothing guards it, because the second cause is
independent of the dial.  The starter owes a release of the slot it was born
holding.

\begin{schema}{StRel0}
  \Delta Dial \\ \Xi Env \\ \Xi Flags \\ \Delta Pc \\ \Xi Dc \\ \Xi RLoc
\where
  pc~aSt = prel0 \\
  slot = slistener \\
  ended' = ended \\
  current' = current \\
  sock' = sock \\
  capN' = capN \\
  hubreq' = hubreq \\
  debt' = debt \\
  slot' = slot \\
  pc' = pc \oplus \{ aSt \mapsto phub \}
\end{schema}

THE DEFECT (round 9, silent-failure finding 1).  The starter does not release
the slot.  \texttt{start\_on} raises on a dial ended between its publication and
its start, but the slot still says \texttt{listener} with no listener behind
it; the Hub end that follows is refused behind that phantom holder, and nobody
settles the dial.  This is the one place the file differs from
\texttt{dial\_settlement.tex}.

\section{The ends}

\begin{schema}{DStartDisp}
  \Xi Dial \\ \Xi Env \\ \Xi Flags \\ \Delta Pc \\ \Delta Dc \\ \Xi RLoc
\where
  pc~aDi = pidle \\
  dcs' = cdisp \\
  pc' = pc \oplus \{ aDi \mapsto pend \}
\end{schema}

\begin{schema}{DStartUnresp}
  \Xi Dial \\ \Xi Env \\ \Xi Flags \\ \Delta Pc \\ \Delta Dc \\ \Xi RLoc
\where
  pc~aDi = pidle \\
  dcs' = cunresp \\
  pc' = pc \oplus \{ aDi \mapsto pend \}
\end{schema}

A discoverer: a write that failed (the display's end), or the keepalive giving
up (unresponsive).

\begin{schema}{RecEnd}
  \Delta Dial \\ \Xi Env \\ \Xi Flags \\ \Delta Pc \\ \Xi Dc \\ \Xi RLoc
\where
  pc~aDi = pend \\
  current = fset \\
  slot' = slot \\
  debt' = debt \\
  current' = current \\
  sock' = sock \\
  capN' = capN \\
  hubreq' = hubreq \\
  ((ended = nocause \land ended' = dcs) \lor (ended \neq nocause \land ended' = ended)) \\
  pc' = pc \oplus \{ aDi \mapsto psu \}
\end{schema}

\begin{schema}{RecGone}
  \Xi Dial \\ \Xi Env \\ \Xi Flags \\ \Delta Pc \\ \Xi Dc \\ \Xi RLoc
\where
  pc~aDi = pend \\
  current = fclear \\
  pc' = pc \oplus \{ aDi \mapsto pdone \}
\end{schema}

The cause is recorded, first wins.  A dial no longer current is left alone.

\begin{schema}{SPStart}
  \Xi Dial \\ \Xi Env \\ \Xi Flags \\ \Delta Pc \\ \Xi Dc \\ \Xi RLoc
\where
  pc~aSp = pidle \\
  pc' = pc \oplus \{ aSp \mapsto phub \}
\end{schema}

The stopper: \texttt{stop} (when a listener exists) or \texttt{close}.  Both
end by asking for the Hub's end.  A \texttt{stop} on a dial with no listener
returns without ending anything; that is a subset of \texttt{close}.

\begin{schema}{HubReq}
  \Delta Dial \\ \Xi Env \\ \Xi Flags \\ \Delta Pc \\ \Xi Dc \\ \Xi RLoc \\
  a? : ACTOR
\where
  a? \in \{ aSt, aSp \} \\
  pc~a? = phub \\
  current = fset \\
  ended' = ended \\
  slot' = slot \\
  debt' = debt \\
  current' = current \\
  sock' = sock \\
  capN' = capN \\
  hubreq' = fset \\
  pc' = pc \oplus \{ a? \mapsto psu \}
\end{schema}

\begin{schema}{HubGone}
  \Xi Dial \\ \Xi Env \\ \Xi Flags \\ \Delta Pc \\ \Xi Dc \\ \Xi RLoc \\
  a? : ACTOR
\where
  a? \in \{ aSt, aSp \} \\
  pc~a? = phub \\
  current = fclear \\
  pc' = pc \oplus \{ a? \mapsto pdone \}
\end{schema}

The Hub's end sets a request.  It records no cause and captures nothing.

\section{The no-reader settle and the release}

\begin{schema}{SUGone}
  \Xi Dial \\ \Xi Env \\ \Xi Flags \\ \Delta Pc \\ \Xi Dc \\ \Xi RLoc \\
  a? : ACTOR
\where
  a? \in \{ aSt, aSp, aDi \} \\
  pc~a? = psu \\
  current = fclear \\
  pc' = pc \oplus \{ a? \mapsto pdone \}
\end{schema}

\begin{schema}{SUGrant}
  \Delta Dial \\ \Xi Env \\ \Xi Flags \\ \Delta Pc \\ \Xi Dc \\ \Xi RLoc \\
  a? : ACTOR
\where
  a? \in \{ aSt, aSp, aDi \} \\
  pc~a? = psu \\
  current = fset \\
  slot = sfree \\
  slot' = sinline \\
  ended' = ended \\
  debt' = debt \\
  current' = current \\
  sock' = sock \\
  capN' = capN \\
  hubreq' = hubreq \\
  pc' = pc \oplus \{ a? \mapsto pset \}
\end{schema}

\begin{schema}{SURefuse}
  \Delta Dial \\ \Xi Env \\ \Xi Flags \\ \Delta Pc \\ \Xi Dc \\ \Xi RLoc \\
  a? : ACTOR
\where
  a? \in \{ aSt, aSp, aDi \} \\
  pc~a? = psu \\
  current = fset \\
  slot \neq sfree \\
  debt' = fset \\
  ended' = ended \\
  slot' = slot \\
  current' = current \\
  sock' = sock \\
  capN' = capN \\
  hubreq' = hubreq \\
  pc' = pc \oplus \{ a? \mapsto pdone \}
\end{schema}

\texttt{settle\_unread}.  A free slot is taken and the caller settles.  A held
slot is left a debt, in the same step as the refusal: the holder cannot release
between the refusal and the debt.

\begin{schema}{SettleCap}
  \Delta Dial \\ \Xi Env \\ \Delta Flags \\ \Delta Pc \\ \Xi Dc \\ \Xi RLoc \\
  a? : ACTOR
\where
  a? \in \{ aSt, aSp, aDi, aIn \} \\
  pc~a? = pset \\
  slot = sinline \\
  current = fset \\
  (ended \neq nocause \lor hubreq = fset \lor eof = fset) \\
  ((ended \neq nocause \land ended' = ended) \lor
    (ended = nocause \land eof = fset \land ended' = cdisp) \lor
    (ended = nocause \land eof = fclear \land hubreq = fset \land ended' = chub)) \\
  sock' = 0 \\
  current' = fclear \\
  capN' = capN + 1 \\
  slot' = slot \\
  debt' = debt \\
  hubreq' = hubreq \\
  ((ended' \in \{ cdisp, cunresp \} \land sock' > 0 \land s2bad' = fset) \lor
    (\lnot (ended' \in \{ cdisp, cunresp \} \land sock' > 0) \land s2bad' = s2bad)) \\
  ((capN \geq 1 \land s3bad' = fset) \lor (capN = 0 \land s3bad' = s3bad)) \\
  s4bad' = s4bad \\
  ((ended = nocause \land ended' = chub \land eof = fset \land s5bad' = fset) \lor
    (\lnot (ended = nocause \land ended' = chub \land eof = fset) \land
      s5bad' = s5bad)) \\
  pc' = pc \oplus \{ a? \mapsto prel \}
\end{schema}

Settle, by a holder of the slot.  One step: drain to the point of blocking,
then read the outcome once.  An end of file is the display's end and wins over
a Hub request that has not yet been decided; a cause already recorded stands.
The ledger is then captured and the dial is not current.  The flags are set
only by a capture that would violate S2, S3 or S5; S4 holds because
\texttt{ended} is written only by the first recorder.

\begin{schema}{SettleGone}
  \Xi Dial \\ \Xi Env \\ \Xi Flags \\ \Delta Pc \\ \Xi Dc \\ \Xi RLoc \\
  a? : ACTOR
\where
  a? \in \{ aSt, aSp, aDi, aIn \} \\
  pc~a? = pset \\
  current = fclear \\
  pc' = pc \oplus \{ a? \mapsto prel \}
\end{schema}

\begin{schema}{RelOwes}
  \Delta Dial \\ \Xi Env \\ \Xi Flags \\ \Delta Pc \\ \Xi Dc \\ \Xi RLoc \\
  a? : ACTOR
\where
  a? \in \{ aSt, aSp, aDi, aIn \} \\
  pc~a? = prel \\
  slot = sinline \\
  debt = fset \\
  debt' = fclear \\
  ended' = ended \\
  slot' = slot \\
  current' = current \\
  sock' = sock \\
  capN' = capN \\
  hubreq' = hubreq \\
  pc' = pc \oplus \{ a? \mapsto pset \}
\end{schema}

\begin{schema}{RelFree}
  \Delta Dial \\ \Xi Env \\ \Xi Flags \\ \Delta Pc \\ \Xi Dc \\ \Xi RLoc \\
  a? : ACTOR
\where
  a? \in \{ aSt, aSp, aDi, aIn \} \\
  pc~a? = prel \\
  slot = sinline \\
  debt = fclear \\
  slot' = sfree \\
  ended' = ended \\
  debt' = debt \\
  current' = current \\
  sock' = sock \\
  capN' = capN \\
  hubreq' = hubreq \\
  pc' = pc \oplus \{ a? \mapsto pdone \}
\end{schema}

The release.  A debt, left by a refusal that came before it, is honoured by
keeping the slot and settling again.  A debt left after the release finds the
slot free.

\section{The inline reader}

\begin{schema}{IAcq}
  \Delta Dial \\ \Xi Env \\ \Xi Flags \\ \Delta Pc \\ \Xi Dc \\ \Xi RLoc
\where
  pc~aIn = pidle \\
  slot = sfree \\
  slot' = sinline \\
  ended' = ended \\
  debt' = debt \\
  current' = current \\
  sock' = sock \\
  capN' = capN \\
  hubreq' = hubreq \\
  pc' = pc \oplus \{ aIn \mapsto iread \}
\end{schema}

\begin{schema}{IReadIdle}
  \Xi Dial \\ \Xi Env \\ \Xi Flags \\ \Delta Pc \\ \Xi Dc \\ \Xi RLoc
\where
  pc~aIn = iread \\
  pc' = pc \oplus \{ aIn \mapsto ilast \}
\end{schema}

\begin{schema}{IReadNotice}
  \Xi Dial \\ \Xi Env \\ \Xi Flags \\ \Delta Pc \\ \Xi Dc \\ \Xi RLoc
\where
  pc~aIn = iread \\
  (ended \neq nocause \lor hubreq = fset \lor eof = fset) \\
  pc' = pc \oplus \{ aIn \mapsto pset \}
\end{schema}

\begin{schema}{ILastSet}
  \Xi Dial \\ \Xi Env \\ \Xi Flags \\ \Delta Pc \\ \Xi Dc \\ \Xi RLoc
\where
  pc~aIn = ilast \\
  current = fset \\
  (ended \neq nocause \lor hubreq = fset) \\
  pc' = pc \oplus \{ aIn \mapsto pset \}
\end{schema}

\begin{schema}{ILastSkip}
  \Xi Dial \\ \Xi Env \\ \Xi Flags \\ \Delta Pc \\ \Xi Dc \\ \Xi RLoc
\where
  pc~aIn = ilast \\
  (current = fclear \lor (ended = nocause \land hubreq = fclear)) \\
  pc' = pc \oplus \{ aIn \mapsto prel \}
\end{schema}

A caller that holds the slot reads, notices an end (or a request) and owes a
settle, or finishes its read; its last check, \texttt{\_settle\_if\_ended}, then
either finds an end and owes a settle or finds none and releases.  Between the
last check and the release, anyone may end the dial: that is the window that
the debt closes.

\section{The listener thread}

\begin{schema}{RNotice}
  \Xi Dial \\ \Xi Env \\ \Xi Flags \\ \Delta Pc \\ \Xi Dc \\ \Xi RLoc
\where
  pc~aRd = rtop \\
  (hubreq = fset \lor ended \neq nocause \lor eof = fset) \\
  pc' = pc \oplus \{ aRd \mapsto pset \}
\end{schema}

\begin{schema}{RSettleCap}
  \Delta Dial \\ \Xi Env \\ \Delta Flags \\ \Delta Pc \\ \Xi Dc \\ \Xi RLoc
\where
  pc~aRd = pset \\
  slot = slistener \\
  current = fset \\
  ((ended \neq nocause \land ended' = ended) \lor
    (ended = nocause \land eof = fset \land ended' = cdisp) \lor
    (ended = nocause \land eof = fclear \land hubreq = fset \land ended' = chub)) \\
  sock' = 0 \\
  current' = fclear \\
  capN' = capN + 1 \\
  slot' = slot \\
  debt' = debt \\
  hubreq' = hubreq \\
  ((ended' \in \{ cdisp, cunresp \} \land sock' > 0 \land s2bad' = fset) \lor
    (\lnot (ended' \in \{ cdisp, cunresp \} \land sock' > 0) \land s2bad' = s2bad)) \\
  ((capN \geq 1 \land s3bad' = fset) \lor (capN = 0 \land s3bad' = s3bad)) \\
  s4bad' = s4bad \\
  ((ended = nocause \land ended' = chub \land eof = fset \land s5bad' = fset) \lor
    (\lnot (ended = nocause \land ended' = chub \land eof = fset) \land
      s5bad' = s5bad)) \\
  pc' = pc \oplus \{ aRd \mapsto rdone \}
\end{schema}

\begin{schema}{RSettleGone}
  \Xi Dial \\ \Xi Env \\ \Xi Flags \\ \Delta Pc \\ \Xi Dc \\ \Xi RLoc
\where
  pc~aRd = pset \\
  current = fclear \\
  pc' = pc \oplus \{ aRd \mapsto rdone \}
\end{schema}

The listener's loop tests, in turn, the Hub's request (its stop flag), a
recorded end, and the display's end of file, and on any of them settles through
the same step that every other holder uses.  It never records a cause of its
own, and it reads none twice.

\section{Verification}

\begin{verbatim}
  fuzz docs/dial_settlement_start_fail_buggy.tex      (clean)
  probcli docs/dial_settlement_start_fail_buggy.tex -model_check \
    -p MAX_OPERATIONS 100000 -p MAX_INITIALISATIONS 100
  probcli docs/dial_settlement_start_fail_buggy.tex -model_check -nodead -bf \
    -p MAX_OPERATIONS 100000 -p MAX_INITIALISATIONS 100 \
    -goal "current = fset & (ended /= nocause or hubreq = fset) &
           !a.(a : ACTOR => pc(a) : {pidle,pdone,rdone,rnone})"
\end{verbatim}

Expected and obtained: a deadlock after 12 steps, 106 states analysed.  The
trace: the starter fails and owes a release it never makes; a discoverer records
an unresponsive end and is refused behind the phantom listener, as are the
stopper and the starter's own Hub end.  The wedge goal is found at depth 5 with
no display end at all: the starter fails, skips the release, asks for the Hub's
end, and is refused.


\end{document}
